\documentclass[aps,pra,twocolumn,superscriptaddress,groupedaddress]{revtex4-2}
\usepackage[utf8]{inputenc}
\usepackage[T1]{fontenc}
\usepackage{amsmath,amssymb,dsfont,booktabs,listings}
\usepackage[dvipsnames]{xcolor}
\usepackage[colorlinks=true,linkcolor=blue,urlcolor=blue,citecolor=blue]{hyperref}

\newcommand{\C}{\mathds{C}}
\newcommand{\one}{\mathds{1}}
\newcommand{\code}[1]{\texttt{#1}}
\DeclareMathOperator{\tr}{tr}
\DeclareMathOperator{\spec}{spec}
\definecolor{antongreen}{HTML}{229922}
\newcommand{\anton}[1]{{\color{antongreen} Anton: #1}}

\lstset{
  language=Python,
  basicstyle=\ttfamily\scriptsize,
  keywordstyle=\color{RoyalBlue},
  commentstyle=\color{gray},
  stringstyle=\color{BrickRed},
  columns=fullflexible,
  keepspaces=true,
  frame=single,
  framerule=0.3pt,
  xleftmargin=2pt,
  xrightmargin=2pt,
  showstringspaces=false
}

\begin{document}

\title{\code{bound-entangled} and \code{schur-weyl}: lightweight Python tools\\
for entanglement theory and Schur--Weyl duality}

\author{Authors \anton{author list}}
\affiliation{Naturwissenschaftlich-Technische Fakult\"{a}t, Universit\"{a}t Siegen,
Walter-Flex-Stra\ss e 3, 57068 Siegen, Germany}

\date{\today}

\begin{abstract}
We describe two small open-source Python packages written while developing
semidefinite hierarchies for convex-roof entanglement measures based on
Schur--Weyl duality~\cite{mainpaper}. \code{bound-entangled} gathers, in a
single tested package, the constructions of bound entangled states that are
otherwise scattered across the literature, each returned as a plain NumPy
density matrix. \code{schur-weyl} implements the representation theory of the
symmetric group needed to work with the Schur--Weyl decomposition of
$(\C^d)^{\otimes k}$: partitions and tableaux, dimension formulas, characters,
Schur polynomials, Kronecker coefficients, isotypic projectors and the Schur
transform, without requiring a full computer algebra system. Both packages
depend only on NumPy and are available on PyPI.
\end{abstract}

\maketitle

%==========================================================================
\section{Introduction}
%==========================================================================

Numerical work in entanglement theory constantly needs two kinds of
ingredients. The first is a set of \emph{test states}. Bound entangled states,
i.e.\ entangled states with a positive partial transpose
(PPT)~\cite{Horodecki1998}, are the standard benchmark for any entanglement
criterion or measure, since the PPT criterion~\cite{Peres1996} fails on them by
definition. Their constructions, however, are spread over two decades of papers,
with differing conventions, parametrizations and occasional typos, and they are
typically re-implemented from scratch in each new project.

The second ingredient is the \emph{representation theory of the symmetric
group}. Schur--Weyl duality~\cite{Fulton1991} decomposes
\begin{equation}\label{eq:SW}
  (\C^d)^{\otimes k} \cong \bigoplus_{\substack{\lambda\vdash k\\ \ell(\lambda)\le d}}
  \mathcal{S}^\lambda_k \otimes \mathcal{U}^\lambda_d,
\end{equation}
and underlies spectrum estimation~\cite{Keyl2001}, Schmidt-rank
tests~\cite{Christandl2005} and, in our case, polynomial optimization over
entanglement spectra~\cite{mainpaper}. Working with Eq.~\eqref{eq:SW} requires
partitions, Young tableaux, characters $\chi^\lambda$, dimensions $f^\lambda$,
Schur polynomials $s_\lambda$ and the isotypic projectors
\begin{equation}\label{eq:isotypic}
  \Pi^\lambda = \frac{f^\lambda}{k!}\sum_{\pi\in S_k}\chi^\lambda(\pi)\,\pi .
\end{equation}
These objects are available in computer algebra systems such as
SageMath~\cite{sagemath}, but importing such a system into a numerical
pipeline (here an SDP built on NumPy and a conic solver) is heavy: a
multi-gigabyte installation, a separate environment and slow imports, for a
handful of combinatorial functions.

We therefore wrote two small packages, \code{bound-entangled} and
\code{schur-weyl}, which we describe here. Both return plain NumPy arrays or
Python integers and tuples, so they can be dropped into any existing numerical
code. Section~\ref{sec:be} presents \code{bound-entangled},
Sec.~\ref{sec:sw} presents \code{schur-weyl}, Sec.~\ref{sec:example} gives a
worked example that combines both, and Sec.~\ref{sec:tests} describes how they
are tested.

%==========================================================================
\section{\code{bound-entangled}}\label{sec:be}
%==========================================================================

\subsection{Purpose and design}

The goal of \code{bound-entangled} is simple: one library containing many
bound entangled states to test against, instead of scattered constructions.
Each state is produced by a factory function returning its density matrix
$\rho$ as a NumPy array. The online documentation~\cite{be-docs} has one page
per state that writes out the construction, the parameter range, what is known
about the state and the original reference. For example,
\begin{lstlisting}
from bound_entangled.c3_otimes_c3 import horodecki
from bound_entangled.cd_otimes_cd import yu_oh

rho = horodecki(0.5)                  # 9x9
rho = yu_oh(full_dim=3, x=0.5, y=0.1)  # 9x9
\end{lstlisting}

Subpackages are named after the Hilbert space the states live in:
\code{c2\_otimes\_c4}, \code{c3\_otimes\_c3}, \code{c4\_otimes\_c4},
\code{c5\_otimes\_c5}, \code{cd\_otimes\_cd} (families in any local dimension
$d$), \code{cm\_otimes\_cn} (parametric constructions on $\C^m\otimes\C^n$) and
\code{multipartite}. Two utilities, \code{ketbra} and \code{upb}, build outer
products and the bound entangled state supported on the orthogonal complement
of an unextendible product basis (UPB)~\cite{Bennett1999UPB}. The only runtime
dependency is NumPy, so the states can be imported into any numerical code
without pulling in a larger framework; \code{toqito}~\cite{toqito} is used only
in the test suite to certify the properties of each state
(Sec.~\ref{sec:tests}). \anton{confirm once the toqito removal is released} A
TypeScript port with the same content is also provided.

\subsection{Content}

Table~\ref{tab:states} lists the bipartite states with fixed dimension. In
addition, the package provides grid states built from graph
edges~\cite{Lockhart2018} and their hypergraph
generalization~\cite{Krebs2024}, the GenTiles2 UPB family on
$\C^m\otimes\C^n$~\cite{DiVincenzo2003}, random PPT and NPT states by rejection
sampling, and three multipartite constructions: the Smolin
state~\cite{Smolin2001}, its $2n$-qubit generalization~\cite{Augusiak2006}
and quasi-diagonal-symmetric states that are PPT across every
bipartition~\cite{Tura2018}.

\begin{table*}[t]
\caption{Bipartite states with fixed dimensions in \code{bound-entangled}
v0.2.2 \anton{bump to the toqito-free release and rerun the numbers}, evaluated at the parameters shown. All are unit-trace, Hermitian,
positive semidefinite and PPT to machine precision (smallest eigenvalues of
$\rho$ and $\rho^{T_B}$ above $-3\times10^{-16}$). The last column is the
realignment (CCNR) norm $\|R(\rho)\|_1$; a value above $1$ certifies
entanglement~\cite{Rudolph2005,Chen2003}. States with $\|R(\rho)\|_1\le 1$ are
not detected by this criterion; their entanglement follows from the range
criterion or the UPB construction.}
\label{tab:states}
\begin{ruledtabular}
\begin{tabular}{llcccl}
Factory$^{a}$ & State & Space & Rank & $\|R(\rho)\|_1$ & Ref. \\
\hline
\code{horodecki(0.5)} & Horodecki $2\times4$ & $2\times4$ & 5 & 0.8469 & \cite{Horodecki1997} \\
\code{horodecki\_2\_by\_d\_generalized(4,0.5)} & Horodecki $2\times d$ & $2\times4$ & 5 & 0.8469 & \cite{Tura2012} \\
\code{horodecki(0.5)} & Horodecki $3\times3$ & $3\times3$ & 7 & 1.0023 & \cite{Horodecki1997} \\
\code{chessboard\_extremal\_PPT()} & Bru\ss--Peres chessboard & $3\times3$ & 4 & 1.1667 & \cite{Bruss2000} \\
\code{cross\_hatch()} & Cross-hatch grid state & $3\times3$ & 4 & 1.1401 & \cite{Lockhart2018} \\
\code{ncomms6297()} & V\'ertesi--Brunner nonlocal & $3\times3$ & 4 & 1.0038 & \cite{Vertesi2014} \\
\code{steering\_state(0.5, 0.5)} & Steerable bound entangled & $3\times3$ & 4 & 1.0726 & \cite{Moroder2014} \\
\code{tiles\_upb()} & Tiles UPB & $3\times3$ & 4 & 1.0874 & \cite{Bennett1999UPB} \\
\code{pyramid\_upb()} & Pyramid UPB & $3\times3$ & 4 & 1.0975 & \cite{Bennett1999UPB} \\
\code{yu\_oh(3, 0.5, 0.1)} & Yu--Oh nonlocal & $3\times3$ & 4 & 1.0329 & \cite{YuOh2017} \\
\code{piani()} & Benatti--Floreanini--Piani & $4\times4$ & 6 & 1.5000 & \cite{Benatti2004} \\
\code{gen\_tiles1(4)} & GenTiles1 UPB & $4\times4$ & 7 & 0.9913 & \cite{DiVincenzo2003} \\
\code{badziag\_private\_singlet(2)} & B\k{a}dzi\k{a}g private singlet & $4\times4$ & 6 & 1.0858 & \cite{Pal2021} \\
\code{orthogonal\_singlet(2)} & PPT singlet $\rho_{F_2}$ & $4\times4$ & 8 & 1.0858 & \cite{Pal2021} \\
\code{sn3\_grid\_state()} & Schmidt number 3 grid state & $5\times5$ & 11 & 1.1428 & \cite{Krebs2024} \\
\end{tabular}
\end{ruledtabular}
\footnotesize{$^{a}$Each factory lives in the subpackage named after its space, e.g.\ \code{c3\_otimes\_c3}; the generalized Horodecki, Yu--Oh, GenTiles1 and singlet families live in \code{cd\_otimes\_cd} and take their arguments as keywords, e.g.\ \code{yu\_oh(full\_dim=3, x=0.5, y=0.1)}.}
\end{table*}

%==========================================================================
\section{\code{schur-weyl}}\label{sec:sw}
%==========================================================================

\subsection{Purpose and design}

\code{schur-weyl} provides the representation-theoretic tools needed to work
with Eq.~\eqref{eq:SW} numerically, without importing SageMath. It is written
from scratch on top of NumPy and the Python standard library, with no
combinatorics or representation-theory dependencies. The wheel is about
$22$~kB and the package imports in about $0.1$~s, essentially the import time
of NumPy itself.

Conventions are kept elementary: partitions are tuples of non-increasing
positive integers, permutations are 0-indexed one-line tuples, tableaux are
tuples of rows, and operators on $(\C^d)^{\otimes k}$ are dense NumPy arrays.
Expensive combinatorial quantities are cached.

\subsection{Content}

Table~\ref{tab:sw} lists the modules. We highlight three features.

\emph{Isotypic projectors.} \code{isotypic\_proj(lam, d)} returns
$\Pi^\lambda$ of Eq.~\eqref{eq:isotypic} as a $d^k\times d^k$ matrix, built
from class sums of the permutation action so that each conjugacy class is
weighted by a single character value. For larger systems,
\code{apply\_isotypic\_proj(tensor, lam, axes)} applies $\Pi^\lambda$ to
chosen axes of a tensor without materializing the matrix. This is exactly the
operation needed for operators of the form $\Pi^\lambda_A\otimes\one_B$ that
act on Alice's copies only.

\emph{Symmetric functions.} Schur polynomials are evaluated with the
Jacobi--Trudi identity, and the change of basis between power sums and Schur
polynomials is given by the characters,
$p_\mu=\sum_{\lambda\vdash k}\chi^\lambda(\mu)\,s_\lambda$, computed with the
Murnaghan--Nakayama rule. Monomial symmetric functions, Kostka matrices and
their inverses give the other standard change of basis.

\emph{Schur--Weyl measure and Schur transform.}
\code{schur\_weyl\_measure(spectrum, k)} returns the weak Schur sampling
distribution $\lambda\mapsto\tr(\Pi^\lambda\rho^{\otimes k})
=f^\lambda s_\lambda(\spec\rho)$~\cite{Keyl2001}, and
\code{schur\_transform\_matrix(d, n)} returns the full Schur transform as a
real orthogonal matrix that block-diagonalizes both actions, using Young's
orthogonal form for the irreducible representations of $S_n$.

\begin{table}[t]
\caption{Modules of \code{schur-weyl} v0.4.0.}
\label{tab:sw}
\begin{ruledtabular}
\begin{tabular}{lp{4.6cm}}
Module & Content \\
\hline
\code{young\_diagrams} & partitions $\lambda\vdash k$ (with length and part bounds), conjugation, hook lengths, dominance order, corners \\
\code{dimensions} & $f^\lambda$ (hook length formula), $\dim\mathcal{U}^\lambda_d=s_\lambda(1,\dots,1)$ (content formula) \\
\code{tableaux} & standard and semistandard tableaux, Kostka numbers and matrices, inverse Kostka matrices \\
\code{symmetric\_group} & composition, inverse, cycle type, conjugacy classes and their sizes \\
\code{character} & $\chi^\lambda(\mu)$ via Murnaghan--Nakayama, character tables \\
\code{symmetric\_functions} & $s_\lambda$ (Jacobi--Trudi), $p_\mu$, $m_\mu$, monomials in the Schur basis \\
\code{kronecker} & Kronecker coefficients $g(\lambda,\mu,\nu,\dots)$ as character sums \\
\code{isotypic} & $\Pi^\lambda$ as a matrix, or applied along tensor axes \\
\code{sw\_measure} & $\tr(\Pi^\lambda\rho^{\otimes k})$ from the spectrum of $\rho$ \\
\code{young\_orthonormal} & Young's orthogonal form, Jucys--Murphy elements \\
\code{schur\_transform} & full Schur transform on $(\C^d)^{\otimes n}$ \\
\end{tabular}
\end{ruledtabular}
\end{table}

%==========================================================================
\section{Worked example}\label{sec:example}
%==========================================================================

The listing below combines both packages to reproduce the first step of the
method of Ref.~\cite{mainpaper}. It takes the $3\times3$ Horodecki state with
$a=0.5$, computes Alice's marginal $\rho_A$, and checks (i) the weak Schur
sampling identity $\tr(\Pi^\lambda\rho_A^{\otimes3})=f^\lambda
s_\lambda(\spec\rho_A)$ for all $\lambda\vdash3$, and (ii) that the linear
entropy $S_L=1-p_2$, lifted to degree $3$ as $p_{(1,1,1)}-p_{(2,1)}$, is
reproduced by the observable $O_P=\sum_\lambda (c_\lambda/f^\lambda)\Pi^\lambda$
with Schur coefficients $c_\lambda$ obtained from the characters.

\begin{lstlisting}
import numpy as np
from bound_entangled.c3_otimes_c3 import horodecki
from schur_weyl import (partitions, dim_specht,
    schur_polynomial, isotypic_proj)
from schur_weyl.character import character

rho = horodecki(0.5)            # 3x3 PPT entangled
rA = np.einsum("ijkj->ik", rho.reshape(3, 3, 3, 3))
spec = list(np.linalg.eigvalsh(rA))
rA3 = np.kron(np.kron(rA, rA), rA)

# (i) tr(Pi^lam rA^{x3}) = f^lam s_lam(spec rA)
for lam in partitions(3):
    lhs = np.trace(isotypic_proj(lam, 3) @ rA3).real
    rhs = dim_specht(lam) * schur_polynomial(lam, spec)
    print(lam, round(lhs, 12), round(rhs, 12))

# (ii) S_L = 1 - p_2 = p_(1,1,1) - p_(2,1) at k = 3
c = {lam: character(lam, (1, 1, 1))
          - character(lam, (2, 1)) for lam in partitions(3)}
O_P = sum(c[l] / dim_specht(l) * isotypic_proj(l, 3)
          for l in partitions(3))
print(c, np.trace(O_P @ rA3).real,
      1 - np.trace(rA @ rA).real)
\end{lstlisting}

The output is
\begin{lstlisting}[language={}]
(3,) 0.376 0.376
(2, 1) 0.588 0.588
(1, 1, 1) 0.036 0.036
{(3,): 0, (2, 1): 2, (1, 1, 1): 2} 0.66 0.66
\end{lstlisting}
(the last two numbers agree to $10^{-16}$). The coefficients
$c_\lambda=(0,2,2)$ give $S_L=2s_{(2,1)}+2s_{(1,1,1)}$ on normalized spectra,
and $O_P\otimes\one_B$ is the objective of the level-3 SDP of
Ref.~\cite{mainpaper} for the linear entropy.

Similarly, applying $\Pi^{(1,1,1)}$ to Alice's three copies of a pure state
with \code{apply\_isotypic\_proj} gives
$\tr\big[(\Pi^{(1,1,1)}_A\otimes\one_B)\,\psi^{\otimes3}\big]$ equal to
$4\times10^{-18}$ for a random state of Schmidt rank $2$ and $1.0\times10^{-2}$
for Schmidt rank $3$, in agreement with the Schmidt-rank test of
Ref.~\cite{Christandl2005}.

%==========================================================================
\section{Testing}\label{sec:tests}
%==========================================================================

Both packages are tested with \code{pytest}, and most tests in
\code{bound-entangled} live in the docstrings. Every state factory is checked
to return a valid density matrix (Hermitian, positive semidefinite, unit
trace) that is PPT and entangled, the latter using the separability checks of
\code{toqito}~\cite{toqito}, which is installed only with the test extra.
Table~\ref{tab:states} reproduces these properties independently.

\code{schur-weyl} is tested against identities with independently known
values: $\sum_{\lambda\vdash k}(f^\lambda)^2=k!$; the completeness relation
$\sum_\lambda\Pi^\lambda=\one$ and idempotence $(\Pi^\lambda)^2=\Pi^\lambda$;
row and column orthogonality of the character tables; the agreement between
$\dim\mathcal{U}^\lambda_d$, $s_\lambda(1,\dots,1)$ and the number of
semistandard tableaux; the identity $\sum_\lambda f^\lambda
s_\lambda(\alpha)=(\sum_i\alpha_i)^k$ behind the Schur--Weyl measure; small
Kronecker coefficients; and the block structure of the Schur transform.
\anton{check this list against the actual test suite}

%==========================================================================
\section{Availability}
%==========================================================================

Both packages require Python $\ge3.9$ and are installed with
\begin{lstlisting}[language=bash]
pip install bound-entangled schur-weyl
\end{lstlisting}
\begin{itemize}
\item \code{bound-entangled} v0.2.2 \anton{new version}, MIT license,
\url{https://github.com/AVlyx/bound-entangled}, documentation at
\url{https://avlyx.github.io/bound-entangled/}.
\item \code{schur-weyl} v0.4.0, MIT license,
\url{https://github.com/AVlyx/schur-weyl}, documentation at
\url{https://avlyx.github.io/schur-weyl/}.
\end{itemize}
The versions described here are archived on Zenodo~\anton{add DOIs}.
Contributions of further states and representation-theoretic tools are
welcome through the GitHub repositories.

{\em AI statement.} \anton{}
An initial draft was made with AI, then improved and partially  rewritten by the authors, and the final version was refined with AI.

{\em Acknowledgements.} \anton{acknowledgements}

\begin{thebibliography}{99}

\bibitem{mainpaper} Authors, \emph{Entanglement of formation and LOCC
conversion from Schur--Weyl duality}, in preparation. \anton{arXiv number}

\bibitem{Horodecki1998} M.~Horodecki, P.~Horodecki, and R.~Horodecki,
Phys.\ Rev.\ Lett.\ \textbf{80}, 5239 (1998).

\bibitem{Peres1996} A.~Peres, Phys.\ Rev.\ Lett.\ \textbf{77}, 1413 (1996).

\bibitem{Fulton1991} W.~Fulton and J.~Harris, \emph{Representation Theory: A
First Course} (Springer, New York, 1991).

\bibitem{Keyl2001} M.~Keyl and R.~F.~Werner, Phys.\ Rev.\ A \textbf{64},
052311 (2001).

\bibitem{Christandl2005} M.~Christandl and G.~Mitchison, Commun.\ Math.\
Phys.\ \textbf{261}, 789 (2006).

\bibitem{sagemath} The Sage Developers, \emph{SageMath, the Sage Mathematics
Software System}, \url{https://www.sagemath.org}.

\bibitem{be-docs} \url{https://avlyx.github.io/bound-entangled/}.

\bibitem{Bennett1999UPB} C.~H.~Bennett, D.~P.~DiVincenzo, T.~Mor, P.~W.~Shor,
J.~A.~Smolin, and B.~M.~Terhal, Phys.\ Rev.\ Lett.\ \textbf{82}, 5385 (1999),
arXiv:quant-ph/9808030.

\bibitem{toqito} V.~Russo, J.~Of.~Open Source Softw.\ \textbf{6}, 3082 (2021).

\bibitem{Lockhart2018} J.~Lockhart, O.~G\"{u}hne, and S.~Severini,
Phys.\ Rev.\ A \textbf{97}, 062340 (2018), arXiv:1705.09261.

\bibitem{Krebs2024} R.~Krebs and M.~Gachechiladze, arXiv:2402.12966.

\bibitem{DiVincenzo2003} D.~P.~DiVincenzo, T.~Mor, P.~W.~Shor, J.~A.~Smolin,
and B.~M.~Terhal, Commun.\ Math.\ Phys.\ \textbf{238}, 379 (2003),
arXiv:quant-ph/9908070.

\bibitem{Smolin2001} J.~A.~Smolin, Phys.\ Rev.\ A \textbf{63}, 032306 (2001),
arXiv:quant-ph/0001001.

\bibitem{Augusiak2006} R.~Augusiak and P.~Horodecki, arXiv:quant-ph/0411142.
\anton{journal ref if published}

\bibitem{Tura2018} J.~Tura, A.~Aloy, R.~Quesada, M.~Lewenstein, and
A.~Sanpera, Quantum \textbf{2}, 45 (2018), arXiv:1706.09423.

\bibitem{Rudolph2005} O.~Rudolph, Quantum Inf.\ Process.\ \textbf{4}, 219
(2005).

\bibitem{Chen2003} K.~Chen and L.-A.~Wu, Quantum Inf.\ Comput.\ \textbf{3},
193 (2003).

\bibitem{Horodecki1997} P.~Horodecki, Phys.\ Lett.\ A \textbf{232}, 333
(1997), arXiv:quant-ph/9703004.

\bibitem{Tura2012} J.~Tura, R.~Augusiak, P.~Hyllus, M.~Ku\'s, J.~Samsonowicz,
and M.~Lewenstein, Phys.\ Rev.\ A \textbf{85}, 060302(R) (2012),
arXiv:1203.3711. \anton{this paper is about four-qubit symmetric PPT states;
check it is the right source for the $2\times d$ Horodecki generalization}

\bibitem{Bruss2000} D.~Bru\ss{} and A.~Peres, Phys.\ Rev.\ A \textbf{61},
030301(R) (2000), arXiv:quant-ph/9911056.

\bibitem{Vertesi2014} T.~V\'ertesi and N.~Brunner, Nat.\ Commun.\ \textbf{5},
5297 (2014).

\bibitem{Moroder2014} T.~Moroder, O.~Gittsovich, M.~Huber, and O.~G\"{u}hne,
Phys.\ Rev.\ Lett.\ \textbf{113}, 050404 (2014), arXiv:1405.0262.

\bibitem{YuOh2017} S.~Yu and C.~H.~Oh, Phys.\ Rev.\ A \textbf{95}, 032111
(2017), arXiv:1509.08991.

\bibitem{Benatti2004} F.~Benatti, R.~Floreanini, and M.~Piani,
arXiv:quant-ph/0411095. \anton{journal ref}

\bibitem{Pal2021} K.~F.~P\'al, G.~T\'oth, E.~Bene, and T.~V\'ertesi,
Phys.\ Rev.\ Research \textbf{3}, 023101 (2021).

\end{thebibliography}

\end{document}